\begin{exercisechunk}
\item \textbf{Most Boolean functions require large circuits.}
\label[exercise]{ex:l14-circuit-counting}
\exerciseprofile{2}{2}{2}{\exproof\quad\excalculation\quad\exinterpretation}{%
Use a counting argument to distinguish universal representation from efficient
representation.}
The Boolean universality theorem constructs a circuit for every Boolean
function, but the theorem does not give a small circuit for every function.
This exercise supplies the counting plan, so the task is to carry out its
consequences rather than invent a circuit encoding.

Fix \(r\) input bits. A circuit has \(s\) numbered gates in topological order.
Each gate chooses one of the sixteen functions
\(\{0,1\}^2\to\{0,1\}\), and each of its two ordered incoming wires may
come from an input or an earlier gate. The circuit output may be any input or
gate value. Let \(N_{r,s}\) be the number of Boolean functions representable
by circuits with at most \(s\) gates.
\begin{enumerate}[(a),beginpenalty=10000]
\item Prove the description-count bound
\[
 N_{r,s}
 \le (s+1)(r+s)\bigl(16(r+s)^2\bigr)^s.
\]
The bound may count several descriptions of the same circuit or function.
Explain why such overcounting does not invalidate the inequality.
\item Define
\[
 s_r=\left\lfloor\frac{2^r}{16r}\right\rfloor.
\]
Show that, for all sufficiently large \(r\),
\[
 \log_2 N_{r,s_r}
 \le 2r+3rs_r
 \le 2r+\frac{3}{16}2^r
 <2^r.
\]
You may use \(r+s_r\le2^r\) and \(4+2r\le3r\) once \(r\) is sufficiently
large.
\item There are \(2^{2^r}\) Boolean functions on \(r\) input bits. Deduce
that some Boolean function requires more than \(s_r\) gates. Show further
that the fraction of Boolean functions representable with at most \(s_r\)
gates tends to zero as \(r\to\infty\).
\item State what the argument establishes about circuit size and what it does
not provide about an explicit hard Boolean function.
\end{enumerate}
\begin{hints}
For part~(a), count descriptions gate by gate, using \(r+s\) as an upper
bound on every predecessor choice, and then choose the output. For part~(b),
take base-two logarithms before substituting \(s_r\). Part~(c) is a
pigeonhole argument.
\end{hints}
\end{exercisechunk}
